\begin{proposition}[The support of relative torsion]\hypertarget{tag-01S1}{}\label{inv:A4.2}
\StatementRecord{inv:A4.2}{proved}{01S1}{\SeriesRef{inv:A4.1}}

The torsion of $\Omega_{X/B}^1$ is supported precisely on
$S_1\cup S_2$. Its cusp modules inject into the free logarithmic
relative modules, and their minimal number of generators exceeds
their rank at the double locus. In degree two, the torsion also
has support on the cusp double locus.
This corrects \cite[Remark~9.21, p.~81]{AC2026}, where
degree-one torsion is asserted on $W\cup S_1\cup S_2$;
the support statement in \cite[Remark~1.3, p.~5]{AC2026}
therefore depends on the degree.
\end{proposition}


\begin{proof}
\begin{enumerate}[label=\arabic*.,leftmargin=*]
\item The double and triple charts in \cref{inv:A4.1} exhaust the
singular strata of $W$. Their injections into free logarithmic
modules prove absence of degree-one cusp torsion. Their degree-two
presentations have a nonzero torsion generator at every general
point of a double curve. At a multiple fibre the finite \'etale
comparison
\[
q^*\Omega^1_{N_j/B}=\Omega^1_{J'_j/B}
\]
is an isomorphism; its summand $R/(s^{m_j-1})\,ds$ is nonzero
for $m_j=3,4$. Faithful flatness descends both the torsion and its
support. There are no other critical fibres, and ordinary relative
differentials are locally free on their complement. Thus the
degree-one support is exactly $S_1\cup S_2$, whereas degree two
also has cusp torsion, as asserted.
\end{enumerate}
\end{proof}